\usemodule[memos][fontsize=16pt,fontstyle=ss]
\usemodule[poriginal]

\starttext

\TitlePage[title=这是文章标题,subtitle=这是副标题,author=\moduleparameter{poriginal}{color},]{}

\Topics{这是篇章目录页}

\Topic{双栏标题框组件测试}

\Subject{基础功能测试}

\subsection{完整双栏标题框（默认）}

\startDualHeaderText[number={No.001},title={这是标题}]
这是一个完整的双栏标题框示例。左侧显示编号，右侧显示标题。
\stopDualHeaderText

\subsection{仅显示编号}

\startDualHeaderText[number={A-123},showtitle=0]
这个框只显示编号，不显示标题。
\stopDualHeaderText

\subsection{仅显示标题}

\startDualHeaderText[title={仅标题示例},shownumber=0]
这个框只显示标题，不显示编号。
\stopDualHeaderText

\subsection{无编号和标题}

\startDualHeaderText[number=,title=,shownumber=0,showtitle=0,toffset=0.5\bodyfontsize]
当编号和标题都为空时，不显示标题栏。需要手动调整 toffset 参数。
\stopDualHeaderText

\Subject{背景样式测试}

\subsection{simple 样式（默认）}

\startDualHeaderText[
  number=, title=, shownumber=0, showtitle=0,
  toffset=0.5\bodyfontsize,
  bodybackground=simple
]
这是默认的简单样式，带完整边框。
\stopDualHeaderText

\startDualHeaderText[
  number=, title=, shownumber=0, showtitle=0,
  toffset=0.5\bodyfontsize,
  bodybackground=simple,
  bodybgcolor=softthemecolor,
]
自定义背景色的简单样式。
\stopDualHeaderText

\subsection{symbolbox 样式}

\startDualHeaderText[
  number={Symbol},
  title={符号框样式},
  bodybackground=symbolbox,
  symbol=\usersymbolparameter{Frame}
]
symbolbox 样式会在右下角显示一个符号。
\stopDualHeaderText

\subsection{charbox 样式}

\startDualHeaderText[
  number={Char},
  title={字符图案样式},
  bodybackground=charbox,
  character={$\star$}
]
charbox 样式会在背景显示旋转45度的字符图案。
\stopDualHeaderText

\subsection{framebox 样式}

\startDualHeaderText[
  number={Frame},
  title={棋盘格样式},
  bodybackground=framebox,
]
framebox 样式会在背景显示棋盘格图案。
\stopDualHeaderText

\subsection{shadowbox 样式}

\startDualHeaderText[
  number={Shadow},
  title={阴影框样式},
  bodybackground=shadowbox,
  shadowcolor=lightthemecolor,
]
shadowbox 样式会在右下角显示阴影效果。
\stopDualHeaderText

\subsection{barbox 样式}

\startDualHeaderText[
  number={Bar},
  title={条形样式},
  bodybackground=barbox,
  leftframe=1,
]
barbox 样式会在左侧显示一个色条。
\stopDualHeaderText

\Subject{边框控制测试}

\subsection{仅左边框}

\startDualHeaderText[
  shownumber=0, showtitle=0,
  leftframe=1, rightframe=0, topframe=0, bottomframe=0,
  toffset=0.5\bodyfontsize,
]
只显示左边框。
\stopDualHeaderText

\subsection{仅左右边框}

\startDualHeaderText[
  shownumber=0, showtitle=0,
  leftframe=1, rightframe=1, topframe=0, bottomframe=0,
  toffset=0.5\bodyfontsize,
]
只显示左右边框。
\stopDualHeaderText

\subsection{仅上边框}

\startDualHeaderText[
  shownumber=0, showtitle=0,
  leftframe=0, rightframe=0, topframe=1, bottomframe=0,
  toffset=0.5\bodyfontsize,
]
只显示上边框。
\stopDualHeaderText

\subsection{无边框}

\startDualHeaderText[
  shownumber=0, showtitle=0,
  leftframe=0, rightframe=0, topframe=0, bottomframe=0,
  toffset=0.5\bodyfontsize,
  bodybgcolor=softthemecolor,
]
不显示任何边框。
\stopDualHeaderText

\Subject{宽度比例测试}

\subsection{1:2 比例}

\startDualHeaderText[
  number={No.},
  title={标题宽度是编号的两倍},
  numberratio=1, titleratio=2
]
标题的宽度是编号的两倍。
\stopDualHeaderText

\subsection{1:3 比例}

\startDualHeaderText[
  number={A},
  title={标题宽度是编号的三倍},
  numberratio=1, titleratio=3
]
标题的宽度是编号的三倍。
\stopDualHeaderText

\Subject{圆角测试}

\subsection{圆角效果}

\startDualHeaderText[
  number={圆角},
  title={圆角效果演示},
  radius=10pt
]
使用 radius 参数设置圆角半径。
\stopDualHeaderText

\subsection{大圆角}

\startDualHeaderText[
  number={大圆角},
  title={大圆角效果},
  radius=15pt
]
更大的圆角半径。
\stopDualHeaderText

\Subject{标题栏边框控制测试}

\subsection{隐藏 number 边框}

\startDualHeaderText[
  number={No.001},
  title={隐藏 number 边框},
  numberframe=0,
  titleframe=1,
]
只隐藏 number 区域的边框。
\stopDualHeaderText

\subsection{隐藏 title 边框}

\startDualHeaderText[
  number={No.002},
  title={隐藏 title 边框},
  numberframe=1,
  titleframe=0,
]
只隐藏 title 区域的边框。
\stopDualHeaderText

\subsection{隐藏所有标题栏边框}

\startDualHeaderText[
  number={No.003},
  title={隐藏所有边框},
  numberframe=0,
  titleframe=0,
]
隐藏 number 和 title 的边框。
\stopDualHeaderText

\subsection{自定义边框颜色}

\startDualHeaderText[
  number={No.004},
  title={自定义边框颜色},
  numberframecolor=red,
  titleframecolor=blue,
  numberframe=1,
  titleframe=1,
]
number 边框红色，title 边框蓝色。
\stopDualHeaderText

\Subject{分隔线控制测试}

\subsection{隐藏分隔线}

\startDualHeaderText[
  number={No.001},
  title={隐藏分隔线},
  headerdivider=0,
]
隐藏 number 和 title 之间的竖线分隔线。
\stopDualHeaderText

\subsection{显示分隔线（默认）}

\startDualHeaderText[
  number={No.002},
  title={显示分隔线},
  headerdivider=1,
]
显示 number 和 title 之间的竖线分隔线。
\stopDualHeaderText

\subsection{圆角时分隔线自动隐藏}

\startDualHeaderText[
  number={No.003},
  title={圆角时分隔线自动隐藏},
  radius=10pt,
  headerdivider=1,
]
当有圆角时，分隔线会自动隐藏（避免视觉冲突）。
\stopDualHeaderText

\Subject{titlepos=1 对齐控制测试}

\subsection{number 左对齐，title 右对齐}

\startDualHeaderText[
  number=01,
  title=标题右对齐,
  titlepos=1,
  numberalign=1,
  titlealign=3,
  numberframe=0, titleframe=0, headerdivider=0,
  leftframe=0, rightframe=0, bottomframe=0, topframe=1,
]
number 左对齐，title 右对齐。
\stopDualHeaderText

\subsection{number 右对齐，title 左对齐}

\startDualHeaderText[
  number=02,
  title=标题左对齐,
  titlepos=1,
  numberalign=3,
  titlealign=1,
  numberframe=0, titleframe=0, headerdivider=0,
  leftframe=0, rightframe=0, bottomframe=0, topframe=1,
]
number 右对齐，title 左对齐。
\stopDualHeaderText

\subsection{两者都居中}

\startDualHeaderText[
  number=03,
  title=标题居中,
  titlepos=1,
  numberalign=2,
  titlealign=2,
  numberframe=0, titleframe=0, headerdivider=0,
  leftframe=0, rightframe=0, bottomframe=0, topframe=1,
]
number 和 title 都居中。
\stopDualHeaderText

\subsection{只有 number 时对齐}

\startDualHeaderText[
  number=04,
  showtitle=0,
  titlepos=1,
  numberalign=3,
  numberframe=0, titleframe=0, headerdivider=0,
  leftframe=0, rightframe=0, bottomframe=0, topframe=1,
]
只有 number 时，number 右对齐。
\stopDualHeaderText

\subsection{只有 title 时对齐}

\startDualHeaderText[
  title=只有标题左对齐,
  shownumber=0,
  titlepos=1,
  titlealign=1,
  numberframe=0, titleframe=0, headerdivider=0,
  leftframe=0, rightframe=0, bottomframe=0, topframe=1,
]
只有 title 时，title 左对齐。
\stopDualHeaderText

\Subject{titlepos=0（顶部标题）测试}

\subsection{完整双栏标题框}

\startDualHeaderText[
  number={No.1},
  title={顶部标题},
  titlepos=0,
]
这是 titlepos=0 的效果，标题显示在顶部，占用正文空间。
\stopDualHeaderText

\subsection{仅标题}

\startDualHeaderText[
  shownumber=0,
  title={仅标题},
  titlepos=0,
]
只显示标题，不显示编号。
\stopDualHeaderText

\subsection{仅编号}

\startDualHeaderText[
  shownumber=1,
  showtitle=0,
  number={No.2},
  titlepos=0,
]
只显示编号，不显示标题。
\stopDualHeaderText

\subsection{标题和编号（左对齐）}

\startDualHeaderText[
  number={No.3},
  title={左对齐标题},
  titlepos=0,
  titlealign=1,
]
标题左对齐，编号右对齐。
\stopDualHeaderText

\subsection{标题和编号（右对齐）}

\startDualHeaderText[
  number={No.4},
  title={右对齐标题},
  titlepos=0,
  titlealign=3,
]
标题右对齐，编号左对齐。
\stopDualHeaderText

\subsection{titlepos=0 + 圆角}

\startDualHeaderText[
  number={No.5},
  title={圆角顶部标题},
  titlepos=0,
  radius=5pt,
]
titlepos=0 与圆角效果的组合。
\stopDualHeaderText

\subsection{titlepos=0 + 阴影}

\startDualHeaderText[
  number={No.6},
  title={阴影顶部标题},
  titlepos=0,
  bodybackground=shadowbox,
]
titlepos=0 与阴影效果的组合。
\stopDualHeaderText

\subsection{titlepos=0 + 隐藏分隔线}

\startDualHeaderText[
  number={No.7},
  title={隐藏分隔线},
  titlepos=0,
  headerdivider=0,
]
titlepos=0 时隐藏分隔线。
\stopDualHeaderText

\Subject{组合效果测试}

\subsection{阴影 + 圆角}

\startDualHeaderText[
  number={组合},
  title={阴影+圆角},
  bodybackground=shadowbox,
  radius=8pt,
  shadowcolor=lightthemecolor,
]
阴影和圆角的组合效果。
\stopDualHeaderText

\subsection{左侧条 + 无标题}

\startDualHeaderText[
  number={组合},
  bodybackground=barbox,
  leftframe=1,
  showtitle=0,
]
左侧色条 + 无标题的组合效果。
\stopDualHeaderText

\subsection{字符图案 + 圆角}

\startDualHeaderText[
  number={组合},
  title={字符图案+圆角},
  bodybackground=charbox,
  character={$\heartsuit$},
  radius=10pt,
]
字符图案和圆角的组合效果。
\stopDualHeaderText

\subsection{仅左边框模拟 blockbox}

\startDualHeaderText[
  number={模拟},
  title={模拟 blockbox},
  leftframe=1,
  rightframe=0, topframe=0, bottomframe=0,
  rulethickness=2pt
]
通过设置仅显示左边框，可以模拟 blockbox 的效果。
\stopDualHeaderText

\subsection{自定义样式示例}

\startDualHeaderText[
  number={? 问题探究},
  title={自定义一个新的样式},
  numberbgcolor=orange,
  numbercolor=white,
  titlecolor=red,
  titlebgcolor=transparentcolor,
  headerbgcolor=transparentcolor,
  bodybgcolor=lightorange,
  numberframe=0, titleframe=0,
  leftframe=0, rightframe=0, topframe=1, bottomframe=0,
  numberratio=1, titleratio=3,
]
自定义一个新的样式
\stopDualHeaderText

\stoptext
